Este informe presenta la resolución del problema de optimización ``AniMaraton'', el cual consiste en maximizar la satisfacción obtenida al ver capítulos de diversos animes bajo restricciones de tiempo y energía. Se implementaron y compararon cuatro enfoques algorítmicos: Fuerza Bruta, Programación Dinámica y dos variantes de heurísticas Greedy. Los resultados muestran que, si bien la Programación Dinámica garantiza la solución óptima en tiempos razonables para el dominio del problema, las heurísticas Greedy ofrecen una alternativa extremadamente eficiente con una pérdida de calidad mínima, especialmente en el caso del algoritmo Greedy basado en prefijos. El análisis experimental detalla el comportamiento de cada algoritmo en términos de tiempo de ejecución, uso de memoria y calidad de la solución.
